\subsection{Résolvantes du transport et conservation de la mémoire}
\label{subsec:ii-memoria-resolvente}

L’histoire prolongée détermine les fenêtres
\(I_m=\{9m+1,\ldots,9m+9\}\), pour \(m\geq0\).
Avec la numérotation de cette section, \(I_m=E_m\) pendant le premier cycle.
Soit \(\sigma_m\in\{-1,1\}\) l’orientation conservée par cette histoire ;
pour \(0\leq m<12\), elle coïncide avec \(\Sigma_{12}(m)\). On définit
\[
 a_m=\sigma_m\sum_{t\in I_m}\mathbf1_{\{2,4,6,8\}}(d_t),
 \qquad |a_m|\leq9.
\]
Le registre et le repère mobile se prolongent par
\[
 z_{m+1}=z_m+a_m\mathbf e_{m\bmod12},\qquad y_m=S^{-m}z_m.
\]
L’état initial \(z_0\) conserve l’histoire antérieure ; il vaut zéro pour
un préfixe initialement vide. Dans cette partie, nous écrivons
\(R=S^{-1}\) pour le transport du repère. La mise à jour
\eqref{eq:eta} est maintenue dans toute fenêtre. L’orientation et les événements
ultérieurs s’obtiennent à partir de l’histoire prolongée : le retour du repère
ne présuppose pas la périodicité de la suite d’événements.

\begin{proposition}[Identité de troncature avec état terminal]
Pour $N\geq1$, soient
\[
Y_N(u)=\sum_{m=0}^N y_mu^m,\qquad
H_{\eta,N}(u)=\sum_{m=0}^{N-1}\eta_mu^m.
\]
On a l’identité polynomiale exacte
\begin{equation}
(I-uR)Y_N(u)=y_0+uH_{\eta,N}(u)-u^{N+1}Ry_N.
\label{mem:identidad-finita}
\end{equation}
\end{proposition}
\begin{proof}
Aux degrés $1\leq m\leq N$, le coefficient du membre de gauche est
$y_m-Ry_{m-1}=\eta_{m-1}$, d’après \eqref{eq:eta}.
Ses coefficients aux degrés zéro et $N+1$ sont $y_0$ et $-Ry_N$,
respectivement. La comparaison des coefficients prouve l’égalité.
\end{proof}

Le terme terminal conserve la mémoire transportée lorsqu’on coupe une histoire.
L’identité permet de comparer une évaluation finie à son prolongement
en maintenant explicite l’état qui reste à la frontière.


\subsubsection{Correspondance avec le développement commun de mémoire}

Le changement de numérotation des fenêtres ne change pas le registre. Dans le
premier cycle, les conventions des deux développements sont reliées par
\[
 E^{(I)}_{m+1}=I_m=E^{(II)}_m,\qquad 0\leq m<12.
\]
Les signes et les événements de chaque fenêtre coïncident après ce
changement d’indice. La convention du transport coïncide également :
\(S\mathbf e_j=\mathbf e_{j+1\bmod12}\) et \(R=S^{-1}\).
Par conséquent, la mise à jour \eqref{eq:eta} est la même identité que
l’équation~\textup{(\HMTIIRefOriginal{mem:actualizacion})} du noyau
commun, avec le même état préalable conservé.

La proposition~\HMTIIRefOriginal{mem:prop-serie} démontre l’identité
de la série génératrice, sa convergence pour \(|u|<1\) et la borne uniforme
de queue~\textup{(\HMTIIRefOriginal{mem:cota})}. L’identité finie
\eqref{mem:identidad-finita} ajoute ici le terme de frontière
\(-u^{N+1}Ry_N\) : elle permet de couper l’histoire sans remplacer son état
terminal par zéro. Le paramètre \(u\) compte, dans les deux cas, des fenêtres
de neuf transitions.

L’élimination de composantes de mémoire conserve l’opérateur, la source
et le lecteur selon la proposition~\HMTIIRefOriginal{mem:prop-schur} et
les équations~\textup{(\HMTIIRefOriginal{mem:schur-fuente})} et
\textup{(\HMTIIRefOriginal{mem:schur-lector})}. L’exemple du cycle
de douze fenêtres~\textup{(\HMTIIRefOriginal{mem:ejemplo-ciclo})}
distingue la compression de la résolvante de la résolvante de la compression.
La composition de trois segments et la transitivité de la réduction
conservent leurs preuves avec les identités
\textup{(\HMTIIRefOriginal{mem:cociclo-tres})} et
\textup{(\HMTIIRefOriginal{mem:schur-transitivo})}.
Ces renvois appartiennent au développement intégral inclus dans ce même
article. Le bilan \eqref{mem:balance} et l’état terminal de
\eqref{mem:identidad-finita} conservent, dans leur présente application,
les accroissements et la frontière du registre ; ils ne déterminent pas à eux seuls
de coefficients analytiques supplémentaires.

